%% kapa: exam=1 round=29 year=2018 period=2 form=A subject=accounting subject_ko=회계학 seq=37 number=77
%% source: Q-Net 감정평가사 기출문제 (cst00309) artlSeq=5202808 file=extracted/1008/5202808/2018년 제29회 감정평가사 1차시험 기출문제/2018년 제29회 감정평가사 1차시험 기출문제/2교시 A형.pdf
%% conversion: kordoc 4.24.51 -> tools/gen1.py (auto); review=proofread-sonnet; flags=none
\begin{question}{77}
다음은 활동기준원가계산을 사용하는 제조기업인 (주)감평의 20x1년도 연간 활동원가 예산자료이다. 20x1년에 회사는 제품 A를 1,000단위 생산하였는데 제품 A의 생산을 위한 활동원가는 ￦830,000으로 집계되었다. 제품 A의 생산을 위해서 20x1년에 80회의 재료이동과 300시간의 직접노동시간이 소요되었다. (주)감평이 제품 A를 생산하는 과정에서 발생한 기계작업시간은?
\begin{framedtext}
연간 활동원가 예산자료 \par
\begin{center}
\small
\begin{tabular}{|c|c|c|r|}
\hline
활동 & 활동원가 & 원가동인 & \multicolumn{1}{c|}{원가동인총수량} \\ \hline
재료이동 & ￦4,000,000 & 이동횟수 & 1,000회 \\ \hline
성형 & ￦3,000,000 & 제품생산량 & 15,000단위 \\ \hline
도색 & ￦1,500,000 & 직접노동시간 & 7,500시간 \\ \hline
조립 & ￦1,000,000 & 기계작업시간 & 2,000시간 \\ \hline
\end{tabular}
\end{center}
\end{framedtext}
\begin{choices}
  \item 400시간
  \item 500시간
  \item 600시간
  \item 700시간
  \item 800시간
\end{choices}
\end{question}
\answer{2}
